\section{基于Stacking的集成}
\label{sec:ensemble-stacking}

前两类方法在成员训练机制中规定了聚合方式。已有线性、近邻、森林或提升树等候选时，也可以将其预测组成新特征，学习这些特征怎样共同解释响应。\term{Stacking}（stacked generalization，堆叠）采用这种两层结构：基学习器产生一级预测，\term{元学习器}（meta-learner）把这些预测映射为最终输出\scite{wolpert1992stacked,breiman1996stacked}。

\subsection{Stacking}
\label{sec:ensemble-stacking-model}

\subsubsection{为什么不能直接使用训练内预测}

设第$m$个成员模型为$h_m$。若直接在完整训练集上拟合$h_m$，再把$h_m(\bm{x}_i)$作为样本$i$的元特征，元学习器看到的是成员对自身训练数据的预测。一个容量很大的成员可能近乎记住训练标签，元学习器便会高估它在新样本上的可靠性。训练内预测由已使用当前标签的成员产生，部署预测则面向未见标签的输入；元学习器需要面对后者的生成条件。

Stacking通常用交叉验证生成\term{折外预测}（out-of-fold prediction, OOF prediction）。把训练样本划分为$K$个互不相交的折$V_1,\ldots,V_K$。对每个成员类型$m$和每个折$k$，只用$D_n\setminus V_k$训练临时模型$h_m^{(-k)}$，再预测$V_k$中的样本。若样本$i$属于$V_{k(i)}$，元特征定义为
\begin{equation}
  z_{i,m}
  =
  h_m^{(-k(i))}(\bm{x}_i).
  \label{eq:ensemble-oof-feature}
\end{equation}
把所有成员的折外预测拼成$\bm{z}_i=(z_{i,1},\ldots,z_{i,M})^{\top}$，再由元学习算法$\mathcal{G}$训练
\[
  g=\mathcal{G}\!\left(\{(\bm{z}_i,y_i)\}_{i=1}^{n}\right).
\]
对每个样本，临时成员的拟合数据排除了它所在的折，其标签只在后续元学习器训练中使用。图\ref{fig:ensemble-stacking-oof}展示了折外矩阵的形成过程。

\begin{figure*}[htbp]
  \centering
  \includegraphics[width=130.048mm]{figures/03-models/01-classic-models/generated/ensemble-stacking-oof/figure.png}
  \caption{三折Stacking的折外特征构造。每一折轮流作为留出折，临时成员只在其余两折训练，并为留出折产生预测；拼接三次预测得到完整折外矩阵$\bm{Z}$，元学习器以$\bm{Z}$为输入并使用对应标签拟合。部署阶段的基成员还需按预先规定的方式重训或保留。}
  \label{fig:ensemble-stacking-oof}
\end{figure*}

分类成员若输出$C$维概率向量，一个成员通常贡献$C$列元特征；二分类时也可以只保留一个类别概率。类别硬判断省去了得分中的排序或置信信息，不同成员的得分还可能采用不同尺度。元训练与部署使用一致的输出形式，正则化再控制尺度及相关预测对组合的影响。

\subsubsection{部署模型与全量训练集重训}

折外矩阵用于训练元学习器，但部署时还需要能预测新输入的基成员。一种做法是在完整训练集上重新拟合每类基学习器$\widehat{h}_m$，预测
\[
  \widehat{y}
  =
  g\!\left(
    \widehat{h}_1(\bm{x}),\ldots,\widehat{h}_M(\bm{x})
  \right).
\]
全量重训利用全部开发数据，但训练规模增加可能改变预测分布，元学习器训练时只看过较小训练集产生的输出。也可保留各折成员，先平均它们对新输入的预测，再送入元学习器；这样沿用折内模型，同时增加了存储与推理次数。两方案的输入分布差异仍由完整验证检验。

部署方案确定后，需要评价基成员与元学习器共同形成的预测流程。基学习器的特征处理、类别编码和超参数搜索必须在相应训练折内完成。若先用全部数据选择特征或编码，再只把模型拟合放进交叉验证，信息仍可能从留出折进入元特征。

\subsection{Blending}
\label{sec:ensemble-blending}

\term{Blending}通过一次留出构造元特征，减少交叉验证的重复拟合。将可用于开发的数据分为基训练集$D_{\mathrm{base}}$和组合训练集$D_{\mathrm{blend}}$；基成员只在$D_{\mathrm{base}}$上拟合，对$D_{\mathrm{blend}}$产生预测，元学习器再以这些预测和对应标签训练。每个基成员只需完成一次用于元特征的拟合，流程比完整Stacking短。

训练数据由两层分用：基成员在$D_{\mathrm{base}}$拟合，元学习器只使用$D_{\mathrm{blend}}$的预测与标签。划分较小时，元学习器估计不稳定；划分较大时，基成员训练信息减少。多次随机划分并平均可以减小单次划分的偶然性，但它重新引入多次训练，已经不再具有单次Blending的全部成本优势。

部署时若把$D_{\mathrm{blend}}$加入基成员重新训练，基预测分布可能相对元学习器训练时发生变化；若不加入，则没有充分利用已有标签。重训方式因此属于模型配置，需要先确定，再用未参与选择的数据评价实际组合。

\subsection{受约束线性Stacking}
\label{sec:ensemble-constrained-linear-stacking}

若希望以较少参数学习成员贡献，可以在回归折外矩阵$\bm{Z}\in\mathbb{R}^{n\times M}$上学习线性组合
\begin{equation}
  (\widehat{b},\widehat{\bm{w}})
  \in
  \argmin_{b,\bm{w}}
  \frac{1}{n}
  \sum_{i=1}^{n}
  \ell\!\left(
    y_i,
    b+\bm{w}^{\top}\bm{z}_i
  \right)
  +\lambda\lVert\bm{w}\rVert_2^2.
  \label{eq:ensemble-linear-stacking}
\end{equation}
权重约束控制组合允许的尺度和方向。在$b=0$时，仅要求$w_m\geq0$得到非负线性组合，仍可能放大预测；仅要求$\sum_mw_m=1$得到仿射组合，仍可因负权重而外推。只有同时满足$b=0$、$w_m\geq0$和$\sum_mw_m=1$，才是成员预测的凸组合，输出必定位于成员预测的凸包内。非负约束排除了成员之间的正负抵消，归一化进一步限制整体尺度；有效校正若需要负权重或截距，约束会缩小候选类并增加逼近误差。

权重大小同时受输出尺度、成员相关性和损失函数影响。两个成员预测近乎相同，线性模型可以在它们之间任意分配相近作用；一个成员输出的数值范围更大，也会改变未标准化正则项下的权重。估计权重反映给定尺度、相关性与损失下的联合拟合，而非成员独立贡献。对同一类别集合上的概率向量，使用同一组满足上述凸组合条件的权重逐类组合，结果保持概率单纯形约束，样本外校准再检查概率与实际频率的对应。

\subsection{Super Learner}
\label{sec:ensemble-super-learner}

\term{Super Learner}把候选学习器库、交叉验证风险与组合权重放进一个明确程序。给定预先规定的候选库和损失，先生成每个候选的折外预测，再在指定组合类中选择交叉验证风险较小的组合；组合类可以是最佳单成员、非负线性组合或其他受控元模型。

在独立同分布、损失有界或满足相应矩条件、候选库与交叉验证方案受控等条件下，Super Learner的理论结果把其风险与候选组合类中的oracle风险联系起来：样本量增加时，交叉验证选择的组合能够渐近接近该库中最优候选或最优组合的风险\scite{vanderlaan2007superlearner}。oracle在预先给定的候选库、组合类和损失中取最优，其范围由这些设计确定。

这一比较依赖三类信息：候选库能提供什么预测，折外损失能多准确地估计目标风险，以及元模型复杂度是否得到样本支持。渐近接近oracle与有限样本中的实际改进需分别检查。若数据存在时间、群组或空间依赖，交叉验证划分还必须与部署时的新样本定义一致。

\subsection{模型对比与总结}
\label{sec:ensemble-stacking-comparison}

表\ref{tab:ensemble-stacking-comparison}比较Stacking家族的元特征来源与组合约束。标准Stacking和Super Learner都依赖折外预测；Blending以一次留出降低训练次数；受约束线性Stacking通过限制权重控制组合层容量。模型名称之外，还应记录预处理、基模型重训和元输出类型，否则无法复现实际预测链路。

\begin{table*}[htbp]
  \centering
  \small
  \begin{tabularx}{\linewidth}{@{}>{\raggedright\arraybackslash}p{25mm}*{5}{>{\raggedright\arraybackslash}X}@{}}
    \toprule
    方法 & 元特征来源 & 数据利用 & 组合规则 & 主要成本 & 主要风险 \\
    \midrule
    Stacking & $K$折折外预测 & 每个样本都产生元特征 & 可线性或非线性 & 每类成员需多次折内拟合 & 折外流程不完整造成泄漏 \\
    Blending & 固定留出集预测 & 元学习器只用留出集 & 可线性或非线性 & 基成员用于元特征时只拟合一次 & 划分敏感与重训分布变化 \\
    受约束线性Stacking & 通常为折外预测 & 与所用划分一致 & 非负、和为一或正则化线性组合 & 组合层较轻量 & 约束可能排除有效校正 \\
    Super Learner & 交叉验证预测与风险 & 由交叉验证方案决定 & 在预定库和组合类中选择 & 候选库越大，训练与选择成本越高 & oracle范围和有限样本被夸大 \\
    \bottomrule
  \end{tabularx}
  \caption{Stacking及其代表性变体。表中“数据利用”描述元学习器训练阶段；部署前是否用完整开发数据重训基成员，还需作为流程的一部分单独规定。}
  \label{tab:ensemble-stacking-comparison}
\end{table*}

多层Stacking把一层折外输出继续交给下一层，每增加一层都要重新保证当前样本没有通过上游拟合泄漏到本层特征。输入相关的门控组合则让权重随$\bm{x}$变化，它与第\ref{sec:ensemble-dynamic-selection}节的动态成员选择相邻，也与神经网络中的混合专家相接。它们增加了表达能力，同时扩大了数据划分、模型存储和调试成本；增加层数需要在相同数据与预算下取得新增预测收益，才能支持相应成本。
